% !Mode:: "TeX:UTF-8"
%# -*- coding:utf-8 -*-

%% 使用 njuthesis 文档类生成南京大学学位论文的示例文档
%%作者：njuhan: https://github.com/njuHan
%% 部分提纲取自njuhan的本科毕业论文

\documentclass[winfonts,bachelor,oneside,nobackinfo]{njuthesis}
%% njuthesis 文档类的可选参数有：
%%   nobackinfo 取消封二页导师签名信息。注意，按照南大的规定，是需要签名页的。
%%   phd/master/bachelor 选择博士/硕士/学士论文

\usepackage{lipsum}
\usepackage{listings}
\usepackage{xcolor}
\lstset{breaklines}%这条命令可以让LaTeX自动将长的代码行换行排版
\lstset{extendedchars=false}%这一条命令可以解决代码跨页时，章节标题，页眉等汉字不显示的问题
\lstset{                        %Settings for listings package.
  language=[ANSI]{C},
  backgroundcolor=\color{white},
  basicstyle=\footnotesize,
  breakatwhitespace=false,
  breaklines=true,
  captionpos=b,
  commentstyle=\color{olive},
  directivestyle=\color{blue},
  extendedchars=false,
  % frame=single,%shadowbox
  framerule=0pt,
  keywordstyle=\color{blue}\bfseries,
  morekeywords={*,define,*,include...},
  numbersep=5pt,
  rulesepcolor=\color{red!20!green!20!blue!20},
  showspaces=false,
  showstringspaces=false,
  showtabs=false,
  stepnumber=2,
  stringstyle=\color{purple},
  tabsize=4,
  title=\lstname
}


%%%%%%%%%%%%%%%%%%%%%%%%%%%%%%%%%%%%%%%%%%%%%%%%%%%%%%%%%%%%%%%%%%%%%%%%%%%%%%%
% 设置《国家图书馆封面》的内容，仅博士论文才需要填写

% 设置论文按照《中国图书资料分类法》的分类编号
\classification{0175.2}
% 论文的密级。需按照GB/T 7156-2003标准进行设置。预定义的值包括：
% - \openlevel，表示公开级：此级别的文献可在国内外发行和交换。
% - \controllevel，表示限制级：此级别的文献内容不涉及国家秘密，但在一定时间内
%   限制其交流和使用范围。
% - \confidentiallevel，表示秘密级：此级别的文献内容涉及一般国家秘密。
% - \clasifiedlevel，表示机密级：此级别的文献内容涉及重要的国家秘密 。
% - \mostconfidentiallevel，表示绝密级：此级别的文献内容涉及最重要的国家秘密。
% 此属性可选，默认为\openlevel，即公开级。
\securitylevel{\controllevel}
% 设置论文按照《国际十进分类法UDC》的分类编号
% 该编号可在下述网址查询：http://www.udcc.org/udcsummary/php/index.php?lang=chi
\udc{004.72}
% 国家图书馆封面上的论文标题第一行，不可换行。此属性可选，默认值为通过\title设置的标题。
\nlctitlea{数据中心}
% 国家图书馆封面上的论文标题第二行，不可换行。此属性可选，默认值为空白。
\nlctitleb{网络模型研究}
% 国家图书馆封面上的论文标题第三行，不可换行。此属性可选，默认值为空白。
\nlctitlec{}
% 导师的单位名称及地址
\supervisorinfo{南京大学计算机科学与技术系~~南京市汉口路22号~~210093}
% 答辩委员会主席
\chairman{张三丰~~教授}
% 第一位评阅人
\reviewera{阳顶天~~教授}
% 第二位评阅人
\reviewerb{张无忌~~副教授}
% 第三位评阅人
\reviewerc{黄裳~~教授}
% 第四位评阅人
\reviewerd{郭靖~~研究员}


%%%%%%%%%%%%%%%%%%%%%%%%%%%%%%%%%%%%%%%%%%%%%%%%%%%%%%%%%%%%%%%%%%%%%%%%%%%%%%%
% 设置论文的中文封面

% 论文标题，不可换行
%\title{南大本科毕业论文\LaTeX 模板}
% 如果论文标题过长，可以分两行，第一行用\titlea{}定义，第二行用\titleb{}定义，将上面的\title{}注释掉
 \titlea{半轻衰变$D^+\to \omega(\phi)e^+\nu_e$的研究}
 \titleb{和弱衰变$J/\psi \to D_s^{(*)-}e^+\nu_e$的寻找}

% 论文作者姓名
\author{作者}
% 论文作者联系电话
\telphone{xxxx}
% 论文作者电子邮件地址
\email{sample@smail.nju.edu.cn}
% 论文作者学生证号
\studentnum{12xxxxxxx}
% 论文作者入学年份（年级）
\grade{2012}
% 导师姓名职称
\supervisor{某~~教授}
% 导师的联系电话
\supervisortelphone{}
% 论文作者的学科与专业方向
\major{计算机科学与技术}
% 论文作者的研究方向
\researchfield{分布式计算}
% 论文作者所在院系的中文名称
\department{计算机科学与技术系}
% 论文作者所在学校或机构的名称。此属性可选，默认值为``南京大学''。
\institute{南京大学}
% 论文的提交日期，需设置年、月、日。
\submitdate{2016年\qquad 月\qquad 日}
% 论文的答辩日期，需设置年、月、日。
\defenddate{2016年\qquad 月\qquad 日}
% 论文的定稿日期，需设置年、月、日。此属性可选，默认值为最后一次编译时的日期，精确到日。
\date{2016年6月8日}

%%%%%%%%%%%%%%%%%%%%%%%%%%%%%%%%%%%%%%%%%%%%%%%%%%%%%%%%%%%%%%%%%%%%%%%%%%%%%%%
% 设置论文的英文封面

% 论文的英文标题，不可换行
\englishtitle{\LaTeX \;  NJU thesis template }
% 论文作者姓名的拼音
\englishauthor{Author}
% 导师姓名职称的英文
\englishsupervisor{ Professor}
% 论文作者学科与专业的英文名
\englishmajor{Computer Science and Technology}
% 论文作者所在院系的英文名称
\englishdepartment{Department of Computer Science and Technology}
% 论文作者所在学校或机构的英文名称。此属性可选，默认值为``Nanjing University''。
\englishinstitute{Nanjing University}
% 论文完成日期的英文形式，它将出现在英文封面下方。需设置年、月、日。日期格式使用美国的日期
% 格式，即``Month day, year''，其中``Month''为月份的英文名全称，首字母大写；``day''为
% 该月中日期的阿拉伯数字表示；``year''为年份的四位阿拉伯数字表示。此属性可选，默认值为最后
% 一次编译时的日期。
\englishdate{May 20th, 2016}

%%%%%%%%%%%%%%%%%%%%%%%%%%%%%%%%%%%%%%%%%%%%%%%%%%%%%%%%%%%%%%%%%%%%%%%%%%%%%%%
% 设置论文的中文摘要

% 设置中文摘要页面的论文标题及副标题的第一行。
% 此属性可选，其默认值为使用|\title|命令所设置的论文标题
% \abstracttitlea{数据中心网络模型研究}
% 设置中文摘要页面的论文标题及副标题的第二行。
% 此属性可选，其默认值为空白
% \abstracttitleb{}

%%%%%%%%%%%%%%%%%%%%%%%%%%%%%%%%%%%%%%%%%%%%%%%%%%%%%%%%%%%%%%%%%%%%%%%%%%%%%%%
% 设置论文的英文摘要

% 设置英文摘要页面的论文标题及副标题的第一行。
% 此属性可选，其默认值为使用|\englishtitle|命令所设置的论文标题
\englishabstracttitlea{englishabstracttitlea}
% 设置英文摘要页面的论文标题及副标题的第二行。
% 此属性可选，其默认值为空白
\englishabstracttitleb{nglishabstracttitleb}

%%%%%%%%%%%%%%%%%%%%%%%%%%%%%%%%%%%%%%%%%%%%%%%%%%%%%%%%%%%%%%%%%%%%%%%%%%%%%%%
\begin{document}

%%%%%%%%%%%%%%%%%%%%%%%%%%%%%%%%%%%%%%%%%%%%%%%%%%%%%%%%%%%%%%%%%%%%%%%%%%%%%%%

% 制作国家图书馆封面（博士学位论文才需要）
%\makenlctitle
% 制作中文封面
\maketitle
% 制作英文封面
\makeenglishtitle


%%%%%%%%%%%%%%%%%%%%%%%%%%%%%%%%%%%%%%%%%%%%%%%%%%%%%%%%%%%%%%%%%%%%%%%%%%%%%%%
% 开始前言部分
\frontmatter

%%%%%%%%%%%%%%%%%%%%%%%%%%%%%%%%%%%%%%%%%%%%%%%%%%%%%%%%%%%%%%%%%%%%%%%%%%%%%%%
% 论文的中文摘要
\begin{abstract}
\lipsum[1-2]

%通过改变链路中子流的个数，分配不同的数据流量给不同的链路。

% 中文关键词。关键词之间用中文全角分号隔开，末尾无标点符号。
\keywords{关键词1 \quad 关键词2 }
\end{abstract}

%%%%%%%%%%%%%%%%%%%%%%%%%%%%%%%%%%%%%%%%%%%%%%%%%%%%%%%%%%%%%%%%%%%%%%%%%%%%%%%
% 论文的英文摘要
\begin{englishabstract}
\lipsum[2]

%Rate adaptation can be implemented by adjusting the number of subflows on each path.

% 英文关键词。关键词之间用英文半角逗号隔开，末尾无符号。
\englishkeywords{keyword1\quad keyword2}
\end{englishabstract}

%%%%%%%%%%%%%%%%%%%%%%%%%%%%%%%%%%%%%%%%%%%%%%%%%%%%%%%%%%%%%%%%%%%%%%%%%%%%%%%
% 论文的前言，应放在目录之前，中英文摘要之后
%
\begin{preface}
\lipsum[1]
\vspace{1cm}
\begin{flushright}
作者\\
20xx年夏于南京大学
\end{flushright}

\end{preface}

%%%%%%%%%%%%%%%%%%%%%%%%%%%%%%%%%%%%%%%%%%%%%%%%%%%%%%%%%%%%%%%%%%%%%%%%%%%%%%%
% 生成论文目录
\tableofcontents

%%%%%%%%%%%%%%%%%%%%%%%%%%%%%%%%%%%%%%%%%%%%%%%%%%%%%%%%%%%%%%%%%%%%%%%%%%%%%%%
% 生成插图清单。如无需插图清单则可注释掉下述语句。
\listoffigures

%%%%%%%%%%%%%%%%%%%%%%%%%%%%%%%%%%%%%%%%%%%%%%%%%%%%%%%%%%%%%%%%%%%%%%%%%%%%%%%
% 生成附表清单。如无需附表清单则可注释掉下述语句。
\listoftables

%%%%%%%%%%%%%%%%%%%%%%%%%%%%%%%%%%%%%%%%%%%%%%%%%%%%%%%%%%%%%%%%%%%%%%%%%%%%%%%
% 开始正文部分
\mainmatter

%%%%%%%%%%%%%%%%%%%%%%%%%%%%%%%%%%%%%%%%%%%%%%%%%%%%%%%%%%%%%%%%%%%%%%%%%%%%%%%
% 学位论文的正文应以《绪论》作为第一章
\chapter{绪论}\label{chapter_introduction}
\section{研究背景}
使用.bib文件管理参考文献引用，引用示例：\cite{BHR12}.\par
\lipsum[1]

\section{MPTCP拥塞控制机制}\label{subsec:mptcp_conges}
\lipsum[1]

\section{研究的目的与意义}
\lipsum[1]




\section{论文结构}

\lipsum[1]
\begin{itemize}
\item 第一章，
\item 第二章，

\end{itemize}



\chapter{相关工作}

\lipsum[1]

\chapter{MPTCP速率分配问题分析}

\lipsum[1]

\section{MPTCP速率分配场景}\label{sec:rate}
\lipsum[1]
\subsection{负载均衡}
\lipsum[1]

\begin{figure}[htbp]
\begin{minipage}[t]{0.5\textwidth}
\centering
\includegraphics[width=0.9\textwidth]{./figure/github.jpg}
\caption{共享拥塞链路}
\label{lab:1}
\end{minipage}%
\begin{minipage}[t]{0.5\textwidth}
\centering
\includegraphics[width=0.9\textwidth]{./figure/github.jpg}
\caption{多跳中继}
\label{lab:2}
\end{minipage}
\end{figure}


\chapter{实验验证}
\lipsum[2]

实验硬件设备如图(\ref{img:1})所示。
\begin{figure}[htbp]
\begin{minipage}[t]{0.5\textwidth}
\centering
\includegraphics[width=0.9\textwidth]{./figure/github.jpg}
\caption{实验硬件设备总览}
\label{img:1}
\end{minipage}
\begin{minipage}[t]{0.5\textwidth}
\centering
\includegraphics[width=0.9\textwidth]{./figure/github.jpg}
\caption{实验测量示意图}
\label{img:2}
\end{minipage}
\end{figure}






\chapter{总结与展望}
\lipsum[1]



%%%%%%%%%%%%%%%%%%%%%%%%%%%%%%%%%%%%%%%%%%%%%%%%%%%%%%%%%%%%%%%%%%%%%%%%%%%%%%%
% 致谢，应放在《结论》之后
\begin{acknowledgement}
%thanks
时光荏苒，在南京大学的本科的学习生活即将结束，四年的时间转瞬即逝，这几年的经历必将成为我人生宝贵的财富。在此论文完成之际，谨向这几年来帮助我的老师和同学表达最衷心的感

\end{acknowledgement}

%%%%%%%%%%%%%%%%%%%%%%%%%%%%%%%%%%%%%%%%%%%%%%%%%%%%%%%%%%%%%%%%%%%%%%%%%%%%%%%




% 参考文献。应放在\backmatter之前。
% 推荐使用BibTeX，若不使用BibTeX时注释掉下面一句。
%\nocite{*}
\bibliography{sample}


% 附录，必须放在参考文献后，backmatter前
\appendix
\chapter{MPTCP内核源代码修改}\label{app:1}
\section{函数mptcp\_v4\_subflows()}
\begin{lstlisting}[language=C]
static void mptcp_v4_subflows(struct sock *meta_sk, const struct mptcp_loc4 *loc, struct mptcp_rem4 *rem)
{
   int i;
   int num;
   printk(KERN_INFO "******** Entering mptcp_v4_subflows ********\n");

   initial_my_global_var();
   switch(my_counter)
   {
		case 1 : num = Fir; break;
		case 2 : num = Sec; break;
		case 3 : num = Thi; break;
		default : num = Fir;
	}

	for (i = 1; i < num; i++)
	{
		printk(KERN_INFO "******** in mptcp_v4_subflows i = %d num = %d********\n",i,num);
		mptcp_init4_subsockets(meta_sk, loc, rem);
	}
	printk(KERN_INFO "******** Leaving mptcp_v4_subflows ********\n");
}
\end{lstlisting}

%%%%%%%%%%%%%%%%%%%%%%%%%%%%%%%%%%%%%%%%%%%%%%%%%%%%%%%%%%%%%%%%%%%%%%%%%%%%%%%
% 书籍附件
\backmatter
%%%%%%%%%%%%%%%%%%%%%%%%%%%%%%%%%%%%%%%%%%%%%%%%%%%%%%%%%%%%%%%%%%%%%%%%%%%%%%%
% 作者简历与科研成果页，应放在backmatter之后
\begin{resume}
% 论文作者身份简介，一句话即可。
\begin{authorinfo}
\noindent 韦小宝，男，汉族，1985年11月出生，江苏省扬州人。
\end{authorinfo}
% 论文作者教育经历列表，按日期从近到远排列，不包括将要申请的学位。
\begin{education}
\item[2007年9月 --- 2010年6月] 南京大学计算机科学与技术系 \hfill 硕士
\item[2003年9月 --- 2007年6月] 南京大学计算机科学与技术系 \hfill 本科
\end{education}
% 论文作者在攻读学位期间所发表的文章的列表，按发表日期从近到远排列。
\begin{publications}
\item Xiaobao Wei, Jinnan Chen, ``Voting-on-Grid Clustering for Secure
  Localization in Wireless Sensor Networks,'' in \textsl{Proc. IEEE International
    Conference on Communications (ICC) 2010}, May. 2010.
\item Xiaobao Wei, Shiba Mao, Jinnan Chen, ``Protecting Source Location Privacy
  in Wireless Sensor Networks with Data Aggregation,'' in \textsl{Proc. 6th
    International Conference on Ubiquitous Intelligence and Computing (UIC)
    2009}, Oct. 2009.
\end{publications}
% 论文作者在攻读学位期间参与的科研课题的列表，按照日期从近到远排列。
\begin{projects}
\item 国家自然科学基金面上项目``无线传感器网络在知识获取过程中的若干安全问题研究''
（课题年限~2010年1月 --- 2012年12月），负责位置相关安全问题的研究。
\item 江苏省知识创新工程重要方向项目下属课题``下一代移动通信安全机制研究''
（课题年限~2010年1月 --- 2010年12月），负责LTE/SAE认证相关的安全问题研究。
\end{projects}
\end{resume}

%%%%%%%%%%%%%%%%%%%%%%%%%%%%%%%%%%%%%%%%%%%%%%%%%%%%%%%%%%%%%%%%%%%%%%%%%%%%%%%
% 生成《学位论文出版授权书》页面，应放在最后一页
%\makelicense

%%%%%%%%%%%%%%%%%%%%%%%%%%%%%%%%%%%%%%%%%%%%%%%%%%%%%%%%%%%%%%%%%%%%%%%%%%%%%%%
\end{document}
